% !TEX root = ../presentation.tex
% =============================================================
%  Subsection — Power transmission limits of an AC OHL
%  Part of ac-transmission; \input from ac-transmission.tex,
%  between ohl-configuration.tex and series-compensation.tex.
%
%  Figures in ./fig (tl_*.pdf), all vector, generated by
%  fig/transmission_limits_figures.py from the exact two-port equations
%  (numpy + matplotlib + scipy.optimize.brentq).  Criteria:
%    thermal    P_th = sqrt(3) U_n I_th, I_th = 3 x 1300 A (3x Bluebird, 75 C)
%    voltage    U_S = 1.0 p.u., unity-pf load at R, U_R >= 0.95 p.u.
%    stability  both terminals 1.0 p.u., delta <= 30 deg
%  Case (a) = min of all three (load at receiving end, St. Clair);
%  case (b) = min of thermal and stability (two stiff systems).
%
%  Reference line: 500 kV, 60 Hz, 3x Bluebird ACSR, flat, 35 ft / 18 in,
%  P_nat = 939 MW, P_th = 3377 MW; series bank at the midpoint.
% =============================================================
\providecommand{\acfig}[1]{ac-transmission/fig/#1}

\intermezzo{Power Transmission Limits}{Thermal, voltage and stability limits of an OHL}

% ---------------------------------------------------------------
\begin{frame}{Three limits on the power transfer of a line}
  \begin{columns}[T,onlytextwidth]
    \begin{column}{0.50\textwidth}
      {\small
      \begin{itemize}
        \setlength{\itemsep}{0.45em}
        \item \textbf{Thermal.} Conductor temperature sets the current;
              independent of length. Binding only on short lines.
        \item \textbf{Voltage drop.} The receiving-end voltage must stay
              in its band under load: $U_R\ge0.95$\,p.u.\ with
              $U_S=1.0$\,p.u.
        \item \textbf{Steady-state stability.} $\delta$ is kept well below
              $90^\circ$ so the system survives the loss of a parallel
              element; $\delta\le30^\circ$ leaves a factor 2 on the
              lossless limit.
        \item The \emph{loadability} is the minimum of the three; against
              length it is the St.~Clair curve.
      \end{itemize}}
    \end{column}
    \begin{column}{0.47\textwidth}
      \setlength{\abovedisplayskip}{4pt}\setlength{\belowdisplayskip}{4pt}
      \[ P_{th} = \sqrt{3}\cdot U_n\cdot I_{th} \]
      \[ U_S = \left|A\cdot\underline{U}_R + B\cdot\underline{I}_R\right|
         \le 1.0\,\text{p.u.} \]
      \[ P_{stab} \approx \frac{U^2}{\left(1-k\right)\cdot X_L}\cdot\sin\delta_{max} \]
      \renewcommand{\arraystretch}{1.1}
      {\footnotesize
      \begin{tabular}{@{}l l r@{}}
        \textbf{Reference line} & & \multicolumn{1}{c}{\textbf{Value}} \\
        \hline
        $I_{th}$, $3\times$ Bluebird at 75\,$^\circ$C & A & 3900 \\
        $P_{th}$ & MW & 3377 \\
        $P_{nat}$ & MW & 939 \\
        \hline
      \end{tabular}}
    \end{column}
  \end{columns}
  \atnote{$A$, $B$ two-port constants of the (compensated) line. The curves
  use the exact two-port power--angle relation for $P_{stab}$; the lossless
  form is shown for orientation.}
\end{frame}

% ---------------------------------------------------------------
\begin{frame}{Loadability of the 500\,kV reference line}
  \begin{columns}[c,onlytextwidth]
    \begin{column}{0.55\textwidth}\centering
      \topoimg[width=\linewidth,height=0.64\textheight,keepaspectratio]{\acfig{tl_loadability.pdf}}
    \end{column}
    \begin{column}{0.43\textwidth}
      {\small
      \begin{itemize}
        \setlength{\itemsep}{0.35em}
        \item \textbf{(a)} load at the receiving end, no reactive support:
              thermal to 25\,mi, then voltage drop; falls through
              $1\cdot P_{nat}$ at 125\,mi.
        \item \textbf{(b)} two stiff systems, both ends at 1.0\,p.u.: thermal
              to 60\,mi, then angular stability; falls through
              $1\cdot P_{nat}$ at 240\,mi.
        \item At 300\,mi both cases give 763\,MW ($0.81\cdot P_{nat}$) ---
              the voltage and the angle criterion happen to coincide; the
              line cannot carry its natural power without compensation.
      \end{itemize}}
    \end{column}
  \end{columns}
  \atnote{Uncompensated line. Voltage criterion with a unity power-factor
  load; a lagging load lowers it further, reactive support at the receiving
  end raises it.}
\end{frame}

% ---------------------------------------------------------------
\begin{frame}{Loadability with series compensation}
  \begin{columns}[c,onlytextwidth]
    \begin{column}{0.55\textwidth}\centering
      \topoimg[width=\linewidth,height=0.56\textheight,keepaspectratio]{\acfig{tl_loadability_k.pdf}}
    \end{column}
    \begin{column}{0.43\textwidth}
      \renewcommand{\arraystretch}{1.15}
      {\footnotesize
      \begin{tabular}{@{}c r r r r@{}}
        \textbf{$k$} & \multicolumn{1}{c}{\textbf{$P_{stab}$}}
                     & \multicolumn{1}{c}{\textbf{$P_{U}$}}
                     & \multicolumn{1}{c}{\textbf{(b)}}
                     & \multicolumn{1}{c}{\textbf{thermal to}} \\
                     & \multicolumn{1}{c}{MW} & \multicolumn{1}{c}{MW}
                     & \multicolumn{1}{c}{$P_{nat}$} & \multicolumn{1}{c}{mi} \\
        \hline
        0   &  763 & 763 & 0.81 &  60 \\
        0.3 & 1027 & 826 & 1.09 &  85 \\
        0.5 & 1311 & 844 & 1.40 & 110 \\
        0.7 & 1708 & 796 & 1.82 & 145 \\
        \hline
      \end{tabular}}\\[0.5em]
      {\small
      \begin{itemize}
        \setlength{\itemsep}{0.3em}
        \item Case (b): the bank acts on the stability limit --- $k=0.5$
              lifts the 300\,mi loadability from 0.81 to $1.40\cdot P_{nat}$.
        \item $P_U$ of case (a) barely moves: the $I\cdot X_C$ rise is in
              quadrature with a unity-pf load, so reactive support at the
              receiving end is needed as well.
      \end{itemize}}
    \end{column}
  \end{columns}
  \atnote{Values at 300\,mi, midpoint bank. Curves show case (b), the
  planning case for a line between two strong systems.}
\end{frame}
